\section*{Introdução}
\addcontentsline{toc}{section}{Introdução}
\label{sec:introducao}

A expansão agropecuária costuma ser analisada à luz do dilema entre conservar a vegetação nativa, promovendo a sustentabilidade ambiental, e converter a terra para explorar usos potencialmente mais rentáveis \cite{Angelsen2010}. Essa representação, contudo, simplifica uma decisão que não depende apenas da aptidão agronômica da terra ou da rentabilidade esperada da produção. A conversão da vegetação é uma decisão privada realizada sob preços, custos de produção, restrições de crédito, expectativas, direitos de propriedade, regras ambientais e mecanismos de implementação dessas regras. Os retornos dos usos alternativos da terra e a qualidade da terra ajudam a explicar as escolhas dos proprietários \cite{LubowskiEtAl2008}, enquanto o acesso ao crédito e a implementação da regulação podem produzir respostas distintas conforme a atividade e o perfil do produtor \cite{JungEtAl2022}.

Isto posto, uma mesma regra formal pode produzir respostas diferentes conforme as condições locais de implementação. A análise de arranjos regulatórios na Suíça mostra a relevância dos recursos das autoridades e dos direitos dos proprietários para essa implementação \cite{ViallonEtAl2019}; na Amazônia brasileira, há evidência de que monitoramento e fiscalização podem conter o desmatamento \cite{AssuncaoEtAl2023}. Sob essa perspectiva, a conversão ou conservação da vegetação pode ser tratada como resultado endógeno da interação entre os retornos econômicos dos usos alternativos da terra e os incentivos, custos e riscos associados ao ambiente institucional \cite{PachecoMeyer2022}.

Nessa toada, mudanças nos preços relativos, na rentabilidade agrícola ou nas condições de acesso aos mercados podem aumentar o retorno esperado da conversão, enquanto fiscalização, sanções e restrições legais podem alterar a resposta dos agentes a esses incentivos. A intensidade dessa resposta também pode variar conforme o contexto econômico e as características territoriais \cite{HardingEtAl2021,ChomitzThomas2003}. Portanto, compreender a evolução do estoque de vegetação requer analisar os retornos econômicos associados ao uso da terra e as instituições sob as quais esses retornos são percebidos e transformados em decisões de conversão ou conservação.

Essa relação, ressalte-se, não autoriza antecipar que qualquer aumento de preço aumentará o desmatamento. Na bacia do Tapajós, \citeonline{Braganca2018} encontra que a elevação do preço relativo de lavouras em relação à carne favoreceu a substituição de pastagens por cultivos e reduziu o desmatamento local. A pressão para abrir novas áreas deve, assim, ser distinguida da reorganização da produção em áreas já convertidas. Da mesma forma, o resultado de uma política de conservação deve ser investigado, e não deduzido de sua classificação como instrumento regulatório ou econômico.

Nesse contexto, Mato Grosso constitui um caso empírico especialmente informativo para a investigação dessas relações. O estado combina produção agrícola e pecuária em larga escala com três biomas --- Amazônia, Cerrado e Pantanal --- que apresentam peculiaridades ambientais e econômicas. Além disso, observam-se, dentro de um mesmo território estadual, diferenças relevantes de vegetação, clima, disponibilidade hídrica, sistemas produtivos e custo de oportunidade da terra. A literatura documenta tanto a expansão e a intensificação agrícola quanto a diversidade das trajetórias de uso da terra no estado \cite{VanWeyEtAl2013,SimoesEtAl2020}.

Sendo assim, essa heterogeneidade permite examinar se retornos econômicos e instituições produzem respostas distintas de conversão entre atividades e biomas. Nesse caso, não basta identificar se uma regra ambiental ou uma alteração na rentabilidade afeta, em média, a conversão. Importa verificar se a intensidade dessas respostas depende das condições econômicas e territoriais sobre as quais os incentivos incidem. A heterogeneidade entre Amazônia, Cerrado e Pantanal interessa à pesquisa como parte dos próprios mecanismos econômicos em análise.

\subsection*{Problema, objetivos e justificativa}

Diante desse quadro, o problema central deste artigo consiste em compreender como incentivos econômicos e institucionais moldam a conversão ou manutenção da vegetação em ambientes ecologicamente e economicamente heterogêneos. O \textbf{problema} pode ser sintetizado no seguinte questionamento:

\medskip
\noindent\emph{Como retornos econômicos e instituições condicionam as decisões de conversão ou manutenção da vegetação nativa em Mato Grosso, e em que medida essa relação varia conforme a capacidade produtiva local, as atividades agropecuárias e os biomas?}
\medskip

Essa pergunta desdobra-se em duas dimensões empiricamente relacionadas: (i) se mudanças nos preços dos produtos agropecuários se traduzem em maior conversão onde a capacidade produtiva para esses usos é mais elevada; e (ii) se as condições institucionais modificam a intensidade dessa resposta. Em conjunto, essas dimensões buscam esclarecer como incentivos privados e instituições interagem na relação entre produção e conservação, oferecendo subsídios para o desenho e a focalização territorial de políticas ambientais.

O \textbf{objetivo geral} deste artigo é investigar como e em que intensidade os determinantes observáveis do retorno econômico esperado e o ambiente institucional se relacionam com as decisões de conversão da vegetação nativa em Mato Grosso. De forma específica, o artigo busca: (i) construir medidas de conversão bruta para agricultura e pastagem; (ii) examinar a resposta dessas medidas a variações de preços combinadas com capacidade produtiva predeterminada; (iii) avaliar o condicionamento institucional dessa resposta, distinguindo associação de identificação causal; e (iv) verificar sua heterogeneidade entre atividades e biomas.

A relevância deste trabalho decorre, em termos teóricos, de tratar o \textbf{meio ambiente} como componente \textbf{endógeno} do processo econômico, cujo estoque é transformado por decisões de produção e pelas instituições que as condicionam. Essa relação envolve uma \textbf{dimensão territorial}, na medida em que incentivos e instituições podem gerar respostas distintas conforme as condições locais, e uma \textbf{dimensão intertemporal}, pois a conversão presente altera estoques ambientais e oportunidades econômicas futuras. No âmbito da tese, este artigo examina a formação desse estoque pela decisão de uso da terra; o capítulo 2, dedicado à H2, examina se a vegetação previamente existente modifica a resposta da produtividade agrícola a choques climáticos. A articulação dos capítulos permitirá discutir a contribuição produtiva do meio ambiente, sem atribuir a este primeiro artigo um teste do circuito completo de retroalimentação. A condição de insumo produtivo é compatível com a existência de serviços ambientais cujos benefícios ultrapassam a propriedade e geram externalidades.

Em termos empíricos e aplicados, compreender essas relações pode orientar decisões de produtores sobre conversão, conservação e manejo da terra e subsidiar governos e organizações responsáveis por fiscalização e alocação de recursos ambientais. A contribuição pretendida está em revelar quando condições institucionais e territoriais alteram a resposta aos incentivos econômicos. Isso não pressupõe a ineficácia dos instrumentos existentes nem a superioridade automática da governança local: a literatura admite diferentes combinações de regras, monitoramento e sanções, cuja efetividade depende do contexto \cite{AnderssonEtAl2014}.

\subsection*{Hipóteses}

Dito isso, o artigo examina a primeira hipótese da tese e seus dois desdobramentos:

\medskip
\noindent\textbf{(H1) Retornos, instituições e conversão.}

% CLAIM: H1-C001
% Maiores retornos econômicos esperados dos usos agropecuários elevam a taxa de conversão de vegetação nativa.
\noindent\emph{A taxa de conversão de vegetação nativa para usos agropecuários aumenta quando se eleva o retorno econômico esperado desses usos.}

\medskip
\noindent\textbf{(H1-a) Preços e capacidade produtiva local.}

\noindent\emph{O efeito de aumentos nos preços dos produtos agropecuários sobre a taxa de conversão de vegetação nativa é maior onde a capacidade produtiva local da terra para esses usos é mais elevada.}

\medskip
\noindent\textbf{(H1-b) Condicionamento institucional.}

\noindent\emph{A resposta da taxa de conversão ao retorno econômico esperado dos usos agropecuários é condicionada pelo ambiente institucional, particularmente pelas regras, pelo monitoramento, pela fiscalização e pelas sanções que alteram os custos e os riscos associados à conversão.}
\medskip

Nesta hipótese, vegetação nativa corresponde à cobertura vegetal natural existente no território, enquanto sua conversão designa a substituição dessa cobertura por usos agropecuários, como lavouras ou pastagens. A taxa de conversão refere-se à parcela da vegetação nativa disponível em determinado território e período que é convertida para esses usos, constituindo o principal resultado de interesse. A ausência de conversão observada tampouco equivale, necessariamente, a uma decisão ativa de conservação ou à integridade ecológica da vegetação remanescente.

O retorno econômico esperado dos usos agropecuários corresponde ao benefício privado esperado da utilização produtiva da terra, considerando receitas e custos associados à atividade. Empiricamente, esse retorno não é diretamente observado nas bases reunidas, podendo ser aproximado por seus determinantes econômicos observáveis. Entre eles, destacam-se os preços dos produtos agropecuários e a capacidade produtiva local da terra. A H1-a expressa a expectativa de que choques favoráveis de preços produzam maior pressão sobre a conversão onde a terra oferece maior capacidade produtiva para os usos beneficiados. O teste dessa previsão será, portanto, um teste de uma implicação observável de H1, condicionado à validade da aproximação; não uma estimativa direta do retorno líquido esperado.

O ambiente institucional compreende as regras formais e os mecanismos de implementação que alteram os incentivos econômicos associados à conversão, incluindo direitos de propriedade, regulação ambiental, monitoramento, fiscalização e sanções. A H1-b não pressupõe que essas instituições determinem isoladamente a conversão, mas que modifiquem a intensidade com que alterações nos retornos econômicos esperados se traduzem em conversão efetiva. Quando uma medida representar, de maneira defensável, maior efetividade da fiscalização, espera-se menor sensibilidade da conversão ao retorno. Contagens de autos ou valores de multas, por si sós, não satisfazem essa condição.

As relações previstas em H1, H1-a e H1-b podem apresentar heterogeneidade entre atividades agropecuárias e biomas. Diferentes atividades apresentam tecnologias, estruturas de custos e retornos distintos, enquanto os biomas diferem em condições biofísicas e produtivas. Essa heterogeneidade será tratada como dimensão empírica da análise, sem constituir uma hipótese autônoma. Resultados nulos, de sinal contrário ou restritos a determinados grupos serão discutidos como limites ou evidência contrária às previsões propostas.

\subsection*{Método e organização do artigo}

As hipóteses serão examinadas com dados em painel para Mato Grosso, tendo 2001--2024 como período candidato, conforme a delimitação já adotada na preparação da qualificação. A análise combinará transições anuais de uso da terra, preços agropecuários, medidas predeterminadas de capacidade produtiva e registros institucionais. As especificações considerarão a ordenação temporal entre incentivos e conversão, efeitos fixos territoriais e temporais e controles compatíveis com a identificação. A endogeneidade da fiscalização, a dependência espacial e a validade temporal das medidas de capacidade produtiva receberão tratamento explícito. A amostra final depende da harmonização territorial e da obtenção das conversões brutas, ainda pendentes.

Além desta introdução, o artigo está organizado em cinco seções. A primeira revisa a literatura sobre retornos, instituições e conversão da vegetação. A segunda apresenta Mato Grosso como contexto empírico, descreve os dados e explicita a construção das variáveis. A terceira desenvolve a estratégia econométrica para examinar as hipóteses do artigo. A quarta estabelece os critérios de apresentação e discussão dos resultados a serem obtidos, com atenção à heterogeneidade e à robustez. Por fim, a quinta reúne as contribuições pretendidas e as limitações do estudo. Nesta versão para qualificação, não são apresentados resultados econométricos.

